\section{The posted $K$-group input and its dependencies}
\label{sec:dependencies}

This section states the external vanishing input used in the
application and follows its arithmetic dependencies explicitly.
Mok--Zheng's Theorem~1, p.~2, states
$K_{4m}(\mathbb Z)=0$ for $1\leq m\leq30$ \cite{MZ}.
Their manuscript is posted and unrefereed. The two comparison
maps in Propositions~\ref{prop:k-component} and
\ref{prop:component-control} carry this statement to
$X^{[-2m]}=0$ at every odd prime. For $p\leq61$ these residues
include every nontrivial even character; for $p>61$ they give
thirty distinct characters. The trivial class character is
already zero by the norm map to $\mathbb Q$.

The cases $m=1,2$ are the refereed $K_4$ theorem of Rognes
and the $K_8$ theorem of Dutour Sikiri\'c--Elbaz-Vincent--Kupers--Martinet,
respectively; see \cite[\S2, pp.~2--3]{MZ} for the exact
identifications. In the remaining range, the odd primes below
$2^{31}$ are handled by Hart--Harvey--Ong's Vandiver computation
\cite[\S1, assertion~(2)]{HHO} and the $K$-theory comparison.
The two-primary computation uses Rognes--Weibel and
\cite[VI, Corollary~9.8]{Weibel}. Our Galois application uses
only odd primes.

For odd $p>2^{31}$, Mok--Zheng use a strict Selmer comparison,
an arithmetic upper bound extracted from Alp\"oge's archive,
and exact Hecke determinants. More precisely, put $e=2m$.
Their Proposition~2, pp.~2--3, compares the $p$-primary
$K_{2e}(\mathbb Z)$ obstruction with the strict Selmer group
of $(\mathbb Q_p/\mathbb Z_p)(-e)$, unramified everywhere and
locally trivial at $2$ and $p$. The upper bound used in
\S\S3--5, pp.~3--8, is the pointwise Selmer bound in the
archive, specialized to dimension $2$, the trivial character,
and coefficient ring $\mathbb Z_p$. Thus the dependency is
through that arithmetic statement in the archive. The
numerical Riemann hypothesis for zeta and Dirichlet $L$-functions
does not replace any term in this displayed Selmer comparison.

For $3\leq m\leq26$, the auxiliary fields are
$\mathbb Q(\sqrt{-3})$ and $\mathbb Q(\sqrt{-7})$.
At each Hecke probe, a nonzero Selmer class forces a nonunit
trace difference on one of the admissible classical or
unitary branches. Multiplying the branch factors at each
probe before taking their greatest common divisor is the
logical step that permits the branch to vary with the probe.
The determinant checks in \S\S6--7, pp.~8--11, eliminate
these possibilities. For $27\leq m\leq30$, the argument uses
$\mathbb Q(i)$ and an additional transfer of the archive
construction, described in \S6.5, p.~10, and \S7, p.~11.
That transfer removes one comparison input and replaces six
substantive uses in the source argument. Its public provenance
statement reports that this modified entry point has not
been fully rebuilt.

Sutherland's verification report \cite{Sutherland} records
kernel checks of the original combined archive, including
several independent kernel implementations. Those checks
concern the original archive whose hash is given there.
They do not by themselves check the additional Gaussian
transfer. No new Lean kernel execution is claimed for the
application in this paper.

\subsection{The joint-representation route}

Pintar's posted theorem, p.~1 of \cite{Pintar}, reaches
$m\leq26$ from the arithmetic construction in the same
archive. Its \S2, p.~2, uses the joint-representation
argument: a positive strict Selmer length produces one
admissible integral representation with residual
semisimplification
\[
                 1\oplus\omega^e\oplus\omega^{-e},
\]
valid simultaneously at every Galois element. The logical
map can be stated directly. An admissible local point gives
such an integral representation. If there are no such
representations, the space of local points is empty.
Evaluating the model upper bound at the identity then has
empty covering set and value zero. The Selmer-to-model
inequality forces strict Selmer length zero. This is the
contrapositive used in Pintar's determinant argument.

This route does not use Mok--Zheng's vanishing theorem as
an input. Both dependency paths do use the archive's
arithmetic construction and the $K$-theory comparison.
Pintar's stated range ends at $m=26$, so it gives no
additional route to the final four components. The version
opened here has public identifier 2611.00829v3 and internal
title-page designation ``Version 5, 23 November 2026''.

\subsection{Soul\'e's torsion bound and its exact specialization}

The comparisons above use no general-number-field
torsion-stabilizer bound from Soul\'e's 2003 Lemma~9.
The reviewed joint-representation proof uses the Selmer
and model inequalities just stated; Mok--Zheng and Pintar
do not cite that lemma in their vanishing arguments.
Here is an explicit calculation fixing the scope of the
bound, without identifying unrelated assertions.

Let $F=\mathbb Q(\sqrt5)$, put $t=(\sqrt5-1)/2$, and set
\[
                 A=\begin{pmatrix}0&-1\\1&t\end{pmatrix}
                    \in\operatorname{SL}_2(O_F).
\]
The identities $t^2+t-1=0$ and $A^2-tA+1=0$ give
$A^5=1$, while $A\ne1$. Its order is $5$, exceeding
$1+\max([F:\mathbb Q],2)=3$.
The valid dimension argument is as follows. Regard a
nonidentity order-$p$ endomorphism of $F^N$ as a rational
endomorphism of dimension $[F:\mathbb Q]N$. Its minimal
polynomial has a primitive $p$th root as a zero and hence
contains the irreducible polynomial $\Phi_p$ of degree
$p-1$. Therefore $p\leq[F:\mathbb Q]N+1$. For
$F=\mathbb Q$ this gives $p\leq N+1$, exactly the bound
used in the rational specialization. The displayed
matrix disproves the stated larger-field bound while
the minimal-polynomial argument proves this specialization.

